% Generated draft: requires mathematical and source review.
\documentclass[11pt]{article}
\usepackage[margin=25mm]{geometry}
\usepackage{fontspec}
\setmainfont{Latin Modern Roman}
\usepackage{amsmath,amssymb,mathtools,unicode-math}
\setmathfont{Latin Modern Math}
\usepackage{enumitem,needspace,xurl,hyperref}
\setlength{\parindent}{0pt}
\setlength{\parskip}{5pt}
\setlength{\emergencystretch}{3em}
\newcommand{\N}{\mathbb{N}}
\newcommand{\Z}{\mathbb{Z}}
\newcommand{\Q}{\mathbb{Q}}
\newcommand{\R}{\mathbb{R}}
\newcommand{\C}{\mathbb{C}}
\newcommand{\F}{\mathbb{F}}
\newcommand{\A}{\mathbb{A}}
\newcommand{\PP}{\mathbb{P}}
\newcommand{\E}{\mathbb{E}}
\newcommand{\eps}{\varepsilon}
\newcommand{\dd}{\,\mathrm{d}}
\begin{document}
\section*{UnsolvedMath ID: 30000850}
\textbf{Source problem code:} catalogue entry
\subsection{Definitions and mathematical statement}
Let \(C\) be a general curve of genus \(g\), and let \(R\) be its rational tautological ring in the sense of the cited report, with pieces denoted \(R^i_{(j)}\). The usual smooth projective connected curve setting is intended, but the base field, ambient cycle ring, equivalence relation on cycles, and precise bigrading are not specified in the supplied extraction.

Write \(p(m)\) for the number of partitions of the nonnegative integer \(m\). For the three-argument notation used here, define
\[
p(i,g+1-i,j)=\#\{(\lambda_1,\ldots,\lambda_{i-j}):
\lambda_1\geq\cdots\geq\lambda_{i-j}\geq1,\quad
\textstyle\sum_a\lambda_a=i,\quad
\lambda_a\leq g+1-i\}.
\]
Thus the parts are positive, their number is exactly \(i-j\), and each is at most \(g+1-i\); the empty partition convention gives one partition of zero into zero parts. Outside feasible integer parameters the count is zero. The source does not state the index range of the ring pieces.

The proposed dimension formulas, for a general curve, are
\[
\dim_{\Q}R=p(g+1),\qquad
\dim_{\Q}R^i_{(j)}=p(i,g+1-i,j).
\]
The record explicitly qualifies these predictions by Polishchuk's conjectural identification of an auxiliary ring \(S\) with \(R\). Its assessment describes the underlying conjecture as saying that a proposed ideal gives all relations in the tautological ring. Neither \(S\), that ideal, nor the comparison map is reproduced here. Accordingly the formulas are retained as the source's conditional dimension predictions, and the conjectural identification is not treated as established. The expressions ``general curve'' and ``tautological ring'' must be made precise from the report before interpreting this as a universal statement on a specified moduli locus or cycle ring.
\subsection{Short English statement}
For a general genus-\(g\) curve, do the total and bigraded dimensions of the intended tautological ring equal the stated partition counts, under Polishchuk's conjectural identification with the auxiliary ring \(S\)? The ring and comparison map are not defined in the extraction.
\subsection{Sources}
\begin{itemize}
\item Local UnsolvedMath registry, permanent ID \(30000850\), source code catalogue entry, ``Partition Formulas for Tautological Rings of General Curves.'' Attribution to Polishchuk is supplied for the conjectural identification; no individual proposer is listed.
\item \emph{Komplexe Algebraische Geometrie} (2008), Oberwolfach Reports, DOI 10.4171/owr/2007/47, pp. 2791--2850: \url{https://doi.org/10.4171/owr/2007/47}.
\item Supplied reference, ``Moonen, Relations between tautological cycles'': \url{https://arxiv.org/abs/0706.3478}.
\item Supplied literature triage, checked 2026-08-21: status ``open.'' It identifies the target as Polishchuk's all-relations conjecture for a general curve, and says later work supplies relations and dimension results while continuing to treat the identification as conjectural. This dated assessment derives from supplied OWR curation; time-sensitive status and cited claims require independent verification.
\end{itemize}
\end{document}
